\chapter{Vector measures and the mass of a current}\label{sec:vector-measure-representation}
The positive representation theorem extends to finite-dimensional vector-valued tests. We give the functional-analytic steps explicitly, including the passage from the comass of a form to the mass norm of a multivector. The source is \cite{Sim18}, Chapter 1, Theorem 5.14, printed pp.~37--38, and Chapter 6, Remark 2.7 and equations (2.8)--(2.10), printed pp.~160--161. Simon initially uses the Euclidean norm; the argument below works for an arbitrary finite-dimensional norm. Completeness of the real spaces $L^1$ and $L^2$, elementary integration, and finite-dimensional norm duality are background inputs.

\begin{lemma}[Representation of a bounded functional on a real Hilbert space]\label{lem:hilbert-functional-representation}
If $\mathcal H$ is a real Hilbert space and $L:\mathcal H\to\mathbb R$ is bounded and linear, there is a unique $z\in\mathcal H$ with $L(x)=\langle x,z\rangle$ for all $x$. Moreover $\|z\|=\|L\|$.
\end{lemma}
\begin{proof}
The case $L=0$ is immediate. Otherwise the closed affine set $A=\{x:L(x)=1\}$ is nonempty and its distance $d$ from zero is positive, since $1\leq\|L\|\|x\|$ on $A$. Choose $x_j\in A$ with $\|x_j\|\to d$. Their midpoints belong to $A$, so the parallelogram identity gives
\[
 \|x_j-x_k\|^2=2\|x_j\|^2+2\|x_k\|^2-4\left\|\frac{x_j+x_k}{2}\right\|^2
 \leq2\|x_j\|^2+2\|x_k\|^2-4d^2\longrightarrow0.
\]
Completeness gives a minimizer $x_0\in A$. For $y\in\ker L$, minimizing $\|x_0+ty\|^2$ in $t$ yields $\langle x_0,y\rangle=0$. As $x-L(x)x_0\in\ker L$, we get $L(x)=\langle x,x_0\rangle/\|x_0\|^2$. Set $z=x_0/\|x_0\|^2$. Testing differences against themselves proves uniqueness. Cauchy--Schwarz and testing $z/\|z\|$ give the norm identity.
\end{proof}

\begin{proposition}[The dual of a sigma-finite real integration space]\label{prop:real-lone-duality}
Let $(X,\mu)$ be a sigma-finite measure space. Every bounded linear functional $L:L^1(\mu;\mathbb R)\to\mathbb R$ has the form
\[
 L(f)=\int_X f\eta\,d\mu,\qquad \eta\in L^\infty(\mu),\qquad \|\eta\|_\infty=\|L\|.
\]
The density $\eta$ is unique up to a null set.
\end{proposition}
\begin{proof}
Choose measurable $E_j\uparrow X$ with $\mu(E_j)<\infty$. Extending by zero gives a bounded restriction of $L$ to $L^2(E_j)$, because $\|f\|_1\leq\mu(E_j)^{1/2}\|f\|_2$. Lemma \ref{lem:hilbert-functional-representation} gives a density $\eta_j\in L^2(E_j)$. Write $C=\|L\|$. If $F=\{|\eta_j|>C+\varepsilon\}\subset E_j$ had positive measure, the $L^2$ test $\operatorname{sign}(\eta_j)1_F$ would imply
\[
 (C+\varepsilon)\mu(F)<\int_F|\eta_j|\,d\mu
 =L(\operatorname{sign}(\eta_j)1_F)\leq C\mu(F),
\]
a contradiction. Thus $|\eta_j|\leq C$ almost everywhere. On $E_i\subset E_j$, all $L^2(E_i)$ tests give the same integral for $\eta_i$ and $\eta_j$; testing their difference proves agreement. Discarding countably many null sets glues the densities to $\eta$ on $X$. Truncating both the domain to $E_j$ and the values of an arbitrary $f\in L^1$ approximates it in $L^1$ by bounded functions supported in finite-measure sets. Continuity proves the formula for every $f$. The norm equality follows from $\|\eta\|_\infty\leq C$ and $|L(f)|\leq\|\eta\|_\infty\|f\|_1$. Uniqueness follows by testing the sign of the difference on finite-measure subsets where its absolute value is bounded below.
\end{proof}

\begin{lemma}[The positive variation functional]\label{lem:positive-vector-variation}
Let $E$ be a finite-dimensional real normed space, $U\subset\mathbb R^N$ open, and $L:C_c(U;E)\to\mathbb R$ linear. Suppose for each compact $K\subset U$ there is $C_K<\infty$ with $|L(F)|\leq C_K\|F\|_\infty$ whenever $\operatorname{spt}F\subset K$. For $g\in C_c(U)$ with $g\geq0$, put
\[
 \lambda(g)=\sup\{|L(F)|:F\in C_c(U;E),\ \|F(x)\|\leq g(x)\text{ for every }x\}.
\]
Then $\lambda$ is additive and nonnegatively homogeneous on nonnegative functions. It extends uniquely to a positive linear functional on $C_c(U)$.
\end{lemma}
\begin{proof}
The supremum is finite by the local bound and the support of $g$. Homogeneity is immediate. For additivity from below, choose almost maximizing tests for $g$ and $h$, reversing each sign so that its $L$-value is nonnegative. Their sum is admissible for $g+h$. For the reverse inequality, split any admissible $F$ as
\[
 F_1=\frac{g}{g+h}F,\qquad F_2=\frac{h}{g+h}F
\]
where $g+h>0$, and set both fields to zero elsewhere. These extensions are continuous: $\|F_1\|\leq g$ and $\|F_2\|\leq h$ force continuity at every zero of $g+h$. Each is admissible for its corresponding function, so $|L(F)|\leq\lambda(g)+\lambda(h)$. This proves additivity. Define $\lambda(f)=\lambda(f^+)-\lambda(f^-)$; equality of two decompositions into differences of nonnegative functions follows by cross-adding and using additivity. Hence this extension is linear and positive.
\end{proof}

\begin{theorem}[Polar representation for finite-dimensional tests]\label{thm:normed-vector-riesz}
Under the hypotheses of Lemma \ref{lem:positive-vector-variation}, there are a Radon measure $\mu$ on $U$ and a measurable $\xi:U\to E^*$ such that
\[
 L(F)=\int_U\langle F,\xi\rangle\,d\mu,\qquad
 \|\xi\|_{E^*}=1\quad\mu\text{-almost everywhere},\qquad
 \lambda(g)=\int_Ug\,d\mu\quad(g\in C_c(U),\ g\geq0).
\]
The measure is determined by $\lambda$, and the field is unique $\mu$-almost everywhere for this measure. See \cite{Sim18}, Theorem 5.14; its Hilbert-space statement is the Euclidean-norm special case.
\end{theorem}
\begin{proof}
Theorem \ref{thm:positive-riesz-representation} applied to $\lambda$ gives $\mu$. In particular
\[
 |L(F)|\leq\lambda(\|F\|)=\int_U\|F\|\,d\mu.
\]
Choose a basis $e_1,\ldots,e_d$ of unit vectors of $E$. Each scalar functional $f\mapsto L(fe_i)$ is bounded for the $L^1(\mu)$ norm. The density theorem \ref{thm:radon-lp-density} extends it uniquely from $C_c(U)$ to $L^1(\mu)$. The measure is sigma-finite, since $U$ has a countable compact exhaustion and $\mu$ is finite on compacts. Proposition \ref{prop:real-lone-duality} gives coordinate densities $\eta_i$. Define $\xi$ by $\langle e_i,\xi\rangle=\eta_i$. Finite linear expansion of $F$ proves the integral formula.

For a fixed $q\in E$, the same domination applied to $fq$ and $L^1$ density give $|\langle q,\xi\rangle|\leq\|q\|$ almost everywhere. Take a countable dense subset of the unit ball and discard the union of the associated null sets. Continuity in $q$ then gives $\|\xi\|_{E^*}\leq1$ almost everywhere. For every nonnegative $g\in C_c(U)$, each admissible test in the definition of $\lambda(g)$ satisfies
\[
 |L(F)|\leq\int_Ug\|\xi\|_{E^*}\,d\mu.
\]
Taking the supremum yields
\[
 \int_Ug\,d\mu=\lambda(g)\leq\int_Ug\|\xi\|_{E^*}\,d\mu\leq\int_Ug\,d\mu.
\]
Consequently $\int g(1-\|\xi\|_{E^*})\,d\mu=0$ for every such $g$. Cutoffs equal to one on a compact exhaustion show that this nonnegative integrand vanishes almost everywhere. Positive Riesz uniqueness determines $\mu$, and scalar $L^1$ density determines each coordinate of $\xi$ uniquely.
\end{proof}

\begin{lemma}[Extending a finite-mass current to continuous forms]\label{lem:current-continuous-extension}
A current $T$ with locally finite mass on $U$ has a unique extension to compactly supported continuous $m$-forms that is bounded for uniform comass on each fixed compact support. For open $W\subset U$, its mass $\mathbf M_W(T)$ is also the supremum of $|T(\omega)|$ over continuous compactly supported forms in $W$ of comass at most one.
\end{lemma}
\begin{proof}
Extend a continuous compactly supported form by zero and convolve its coefficients with smooth mollifiers of radii tending to zero. The resulting smooth forms converge uniformly; all sufficiently late supports lie in a common relatively compact open $V\subset U$. The mass bound makes their $T$-values Cauchy. The same estimate applied to the difference of any two approximating sequences proves independence. It also proves continuity and uniqueness of the extension. If the original support lies in $W$, choose the mollifier radii so the approximating supports also lie in $W$. Uniform approximation in the finite-dimensional comass norm gives upper bounds $1+\varepsilon_j$ for the approximating comasses; divide by $1+\varepsilon_j$ to obtain admissible smooth tests. Thus allowing continuous tests does not change the supremum. This is the extension preceding \cite{Sim18}, equation (2.8), with comass in place of the Euclidean norm.
\end{proof}

\begin{proposition}[Open-set mass from the positive variation functional]\label{prop:current-open-mass-variation}
Apply Theorem \ref{thm:normed-vector-riesz} to the continuous extension of a locally finite-mass $m$-current, with $E=\Lambda^m(\mathbb R^N)^*$ carrying comass. Its measure satisfies $\mu(W)=\mathbf M_W(T)$ for every open $W\subset U$. The dual field has mass norm one almost everywhere and need not be simple.
\end{proposition}
\begin{proof}
For a continuous test supported in $W$ with comass at most one, the polar representation gives $|T(\omega)|\leq\mu(W)$. Conversely, if $0\leq g\leq1$ has compact support in $W$, every test for $\lambda(g)$ is such a continuous mass test. Hence $\mathbf M_W(T)\geq\lambda(g)=\int g\,d\mu$. Taking the supremum over these scalar cutoffs gives $\mu(W)$ by Radon regularity, including when it is infinite. The mass norm is precisely the dual of comass, so the unit-norm assertion follows from the vector representation theorem. Nothing in that theorem imposes decomposability of the polar multivector.
\end{proof}
